\scanpage{4}
\diagram{\begin{tikzcd}[ampersand replacement=\&,column sep=small,row sep=small]
A\times A \arrow[d,"\pi"] \arrow[dr,phantom,"{\text{cart}}"{description,draw,ellipse,inner sep=1pt}] \& A\times_S A\times_S B \arrow[l,"\mathrm{pr}_{12}"'] \arrow[d,"{\pi\times \mathrm{id}_B}"] \arrow[dr,"\mathrm{pr}_3"] \& {}\\
A \& A\times_S B \arrow[l,"\mathrm{pr}_1"] \arrow[r,-,"\mathrm{pr}_2"'] \& B
\end{tikzcd}}

\diagram{\begin{tikzcd}[ampersand replacement=\&,column sep=small,row sep=small]
{} \& A\times_S A\times_S B \arrow[r,"\mathrm{pr}_3"] \arrow[d,"{\mathrm{id}_{A\times A}\times\Delta_B=f}"'] \arrow[ddl,bend right=39,"\mathrm{pr}_{12}"'] \arrow[dr,phantom,"{\text{cart}}"{description,draw,ellipse,inner sep=1pt}] \& B \arrow[d,"\Delta_B"]\\
{} \& A\times_S A\times_S B\times_S B \arrow[r,"\mathrm{pr}_{34}"'] \arrow[dl,"\mathrm{pr}'_{12}"] \& B\times_S B\\
A\times_S A \& {} \& {}
\end{tikzcd}}

\[
\begin{aligned}
\Phi(F\overset{\mathbb L}{*} G)&\overset{\text{déf}}{\simeq}\deleted{\unsure{\mathbb{R} \mathrm{pr}_{2*}(\mathbb{L} \mathrm{pr}_1^*(F\overset{\mathbb L}{*} G)\otimes L)}}\inserted{\mathbb{R} \mathrm{pr}_{2*}(\mathbb{L} \mathrm{pr}_1^*(F\overset{\mathbb L}{*} G)\otimes L)}.\\
&\overset{(*)}{=}\mathbb{R} \mathrm{pr}_{2*}\bigl(\mathbb{L} \mathrm{pr}_1^*(\pi_*(F\overset{\mathbb L}{\otimes} G))\inserted{\otimes L}\bigr)\\
&\overset{\text{chgt. de base}}{\simeq}\mathbb{R} \mathrm{pr}_{2*}\bigl(R(\pi\times \mathrm{id}_B)_*(\mathbb{L} \mathrm{pr}_{12}^*(F\overset{\mathbb L}{\otimes}_S G))\otimes L\bigr)\\
&\overset{\text{formule de proj.}}{\simeq}\mathbb{R} \mathrm{pr}_{2*}\bigl(\mathbb{R}(\pi\times \mathrm{id}_B)_*(\mathbb{L} \mathrm{pr}_{12}^*(F\overset{\mathbb L}{\otimes}_S G)\otimes(\pi\times \mathrm{id}_B)^*(L))\bigr)\\
&\overset{\text{transit.}}{\simeq}R \mathrm{pr}_{3*}\bigl(\mathbb{L} \mathrm{pr}_{12}^*(F\overset{\mathbb L}{\otimes}_S G)\otimes(\pi\times \mathrm{id}_B)^*(L)\bigr)
\end{aligned}
\]

Or on a (th. du carré) un isom can.
\[
(\pi\times \mathrm{id}_B)^*(L)\simeq \mathrm{pr}_{13}^*(L)\,\mathrm{pr}_{23}^*(L)
\]
où $\mathrm{pr}_{13},\mathrm{pr}_{23}\colon A\times A\times B\unsure{\mathrel{-}|\to} A\times B$ sont respectivmt $\mathrm{pr}_1\times \mathrm{id}_B$ et $\mathrm{pr}_2\times \mathrm{id}_B$.

On en conclut
\[
\mathbb{L} \mathrm{pr}_{12}^*(F\overset{\mathbb L}{\otimes}_S G)\otimes(\pi\times \mathrm{id}_B)^*(L)\simeq \mathbb{L} f^*\bigl(\underline{\tikz[remember picture,baseline]{\coordinate (quatretermes) at (0,0);}\mathrm{pr}_{12}^{\prime *}(F\overset{\mathbb L}{\otimes}_S G)\otimes \mathrm{pr}_{13}^{\prime *}(L)\otimes \mathrm{pr}_{24}^{\prime *}(L)}\bigr)
\]
où $\mathrm{pr}'_{13},\mathrm{pr}'_{24}$ sont les deux projections $A\times_S A\times_S B\times_S B\longrightarrow A\times_S B$. On a d'ailleurs
\[
\mathrm{pr}_{12}^{\prime *}(F\overset{\mathbb L}{\otimes}_S G)\otimes \mathrm{pr}_{13}^{\prime *}(L)\otimes \mathrm{pr}_{24}^{\prime *}(L)
\overset{\text{ass. de }\unsure{\overset{\mathbb L}{\otimes}_S|\otimes_S|\otimes}}{\simeq}
\tikz[remember picture,baseline]{\coordinate (quatrehaut) at (0,0);}\,\mathrm{pr}_{13}^{\prime *}(F\otimes L)\overset{\mathbb L}{\otimes} \mathrm{pr}_{24}^{\prime *}(G\otimes L)
\]
\deleted{conclut} et on conclut
\[
\begin{aligned}
\Phi(F\overset{\mathbb L}{*} G)&\simeq R \mathrm{pr}_{3*}\bigl(\mathbb{L} f^*(\tikz[remember picture,baseline]{\coordinate (quatrebas) at (0,0);}\hspace{1em})\bigr)\\
&\overset{\text{chgt. de base}}{\simeq}\mathbb{L}\Delta_B^*\bigl(\mathbb{R} \mathrm{pr}_{34*}(\mathrm{pr}_{13}^*(F\otimes_S L)\overset{\mathbb L}{\otimes} \mathrm{pr}_{24}^{\prime *}(G\otimes_S L))\bigr)\\
&\overset{\text{Künneth}}{\simeq}\mathbb{L}\Delta_B^*\bigl(\mathbb{R} \mathrm{pr}_2(F\otimes L)\overset{\mathbb L}{\otimes}_S R \mathrm{pr}_{2*}(G\otimes L)\bigr)\\
&\overset{\text{déf}}{\simeq}\mathbb{L}\Delta_B^*\bigl(\Phi(\unsure{L|F})\overset{\mathbb L}{\otimes}_S\Phi(G)\bigr)\\
&\simeq\Phi(F)\overset{\mathbb L}{\otimes}\Phi(G)
\end{aligned}
\]
\tikz[remember picture,overlay]{\draw (quatrebas) |- (quatrehaut); \draw[->] (quatrehaut) to[bend right=12] (quatretermes);}
cqfd.

\deleted{\uline{Proposition 2}} \inserted{\uline{Théorème}} Soient $A$, \deleted{deux} \inserted{$B$} schémas abéliens sur $S$\inserted{, de dim relative $n$}, \deleted{$\unsure{B|L}$} $L$ un Module inversible birigidifié sur $A\times B$, définissant une \uline{dualité} entre $A$ et $B$. \deleted{Considérons} \inserted{Alors le composé} des deux
\diagram{\begin{tikzcd}[ampersand replacement=\&]
K(A)\arrow[r,"{\unsure{\varphi|\Phi}_L}"] \& K(B)\arrow[r,"{\varphi_{{}^{s}L}}"] \& K(A)
\end{tikzcd}}
\deleted{\unsure{sont|est} \unsure{l'identité}} \inserted{$(-1)^n s_A^*$ (où $s_A$ est la symétrie de $A$)} (donc, par symétrie, on voit que $\Phi_L$ et $\inserted{(-1)^n}s_A^*\Phi_{{}^{s}L}$ sont des isomorphismes inverses l'un de l'autre).
% archive sheet number: 3
